\subsection{Сложение движений}
%%%%%%%%%%%%%%%%%%%%%%%%%%%%%%%%%%%%%%%%%%%%%%%%%%%%
%%%%%%%%%%%%%%%%%%%%%%%%%%%%%%%%%%%%%%%%%%%%%%%%%%%%


 \z[{\begin{wrapfigure}[6]{r}{0.2\textwidth}%
		\vspace{-15pt}%
		\begin{flushright}%
 			\includegraphics[width=0.18\textwidth]{3.4.1}%
		\end{flushright}%
	\end{wrapfigure}%
	}]%
	{Плоский обруч движется так, что в некоторый момент времени скорости концов диаметра $АВ$ лежат в плоскости обруча и перпендикулярны этому диаметру (см. рисунок). Скорость точки $А$ равна $v_A$, а скорость точки $B$ равна $v_B$. Определите скорости концов диаметра $CD$, перпендикулярного $АВ$, в этот же момент времени, считая, что эти скорости также лежат в плоскости обруча.}
%
{$v_c=v_d=\sqrt{\frac{v_a^2+v_b^2}{2}}$}
%%%%%%%%%%%%%%%%%%%%%%%%%%%%%%%%%%%%%%%%%%%%%%%%%%%%
\z[{\begin{wrapfigure}[4]{r}{0.2\textwidth}%
		\vspace{-25pt}%
		\begin{flushright}%
 			\includegraphics[width=0.2\textwidth]{3.4.2}%
		\end{flushright}%
	\end{wrapfigure}%
	}]%
	{Колесо катится без проскальзывания по горизонтальной плоскости. Ускорение центра колеса равно $a_0$. Найдите значения ускорений точек $A$ и $B$ колеса в момент времени, когда ускорение точки $C$ становится равным по модулю $a_0$.}
%
{$a_a=a_b=\sqrt{5}a_0$}
%%%%%%%%%%%%%%%%%%%%%%%%%%%%%%%%%%%%%%%%%%%%%%%%%%%%	
\z{Стержень длиной $2L$ движется в горизонтальной плоскости таким образом, что в некоторый момент времени скорость одного конца стержня равна $v_1$ и направлена под углом $\alpha$ к стержню, скорость второго конца $v_2$. Определите угловую скорость $\omega$ вращения стержня относительно его центра.}
%
{$\omega = \frac{v_1\sin\alpha\pm\sqrt{v_2^2-v_1^2\cos^2\alpha}}{2L}$}
%%%%%%%%%%%%%%%%%%%%%%%%%%%%%%%%%%%%%%%%%%%%%%%%%%%%
 \z[{\begin{wrapfigure}[5]{r}{0.2\textwidth}%
		\vspace{-15pt}%
		\begin{flushright}%
 			\includegraphics[width=0.2\textwidth]{3.4.4}%
		\end{flushright}%
	\end{wrapfigure}%
	}]%
	{Катушка с намотанной на нее нитью лежит на горизонтальной поверхности стола и катится по ней без скольжения под действием нити (см. рисунок). С какой скоростью $v$ будет перемещаться ось катушки, если конец нити тянуть в горизонтальном направлении со скоростью $u$? Радиус внутренней части катушки $r$, внешней -- $R$. Найдите скорость $v_A$ и ускорение $a_A$ точки $А$?}
 %
{$v=u\frac{R}{R-r};~v_a=2u\frac{R}{R-r};~a_A=\frac{u^2R}{(R-r)^2}$}
 
 \z{Решить предыдущую задачу, если нитку тянут под углом $\alpha$ к горизонту.}
%
{$v=u\frac{R}{R\cos{\alpha}-r};~v_a=2u\frac{R}{R\cos{\alpha}-r};~a_A=\frac{u^2R}{(R\cos{\alpha}-r)^2}$}

%%%%%%%%%%%%%%%%%%%%%%%%%%%%%%%%%%%%%%%%%%%%%%%%%%%%